\section{Introduction}
\label{sec:intro}

A fleet of machine-learning accelerators is a large capital asset that behaves
like a queueing system. The chips are bought in bulk, wired into a fixed
hierarchy, and handed out --- minute by minute, for months --- by one program:
the scheduler, which decides which queued job runs now and which waits.

That decision has more leverage than it looks. On a 2{,}048-chip fleet --- the
shipped \texttt{examples/01\_minimal} scenario, not one of the case studies in
\S\ref{sec:cases} --- replacing strict first-in-first-out (FIFO) queueing with a
preempting tiered-priority policy cuts the mean number of waiting evaluation
jobs from about 11 to about 1.3, and their 99th-percentile queue wait from about
5\,h to about 31\,min. Same hardware, same arrivals, same seed; a different
ordering rule. (Those are the figures the example's own README records, to the
precision it records them.)

You cannot run that experiment on the production fleet: an A/B test costs weeks
of accelerator time, and traces record what one scheduler did, never what another
would have done.

\texttt{fleetsim} is an open-source discrete-event simulator that answers such
questions at \emph{fleet altitude}. Altitude here means what one simulated run
represents. One \texttt{fleetsim} run is weeks of a whole fleet: thousands of
jobs arriving, queueing, running, being preempted and occasionally dying with
their nodes. A network simulator's run, by contrast, is one already-placed job
executing a single training step. Given a fleet in YAML, arriving jobs and a
failure model, \texttt{fleetsim} reports the occupancy, queue waits, preemption
rates and goodput a policy delivers --- four quantities
\S\ref{sec:threenumbers} and \S\ref{sec:failures} define, and whose differences
turn out to matter more than any one of them alone. It models no packets,
collectives or kernels --- per-job speed is one analytical multiplier ---
because its questions concern thousands of jobs contending for one machine, not
one job's step time.

We claim three contributions, in ascending order of importance.

\begin{enumerate}
\item \textbf{A fleet-altitude model} (\S\ref{sec:models}): one capacity tree
from metro down to node, spanning NVIDIA-style clusters and TPU-style pods, with
gang scheduling, locality constraints, preemption cost, quota, failures and
drains in one engine; the simulators we surveyed each model a subset
(\S\ref{sec:related}).
\item \textbf{A generative traffic model} (\S\ref{sec:traffic}): arrivals, job
sizes and durations from published cluster characterizations rather than a single
Poisson assumption, each parameter recorded with its source --- including those
that are order-of-magnitude priors rather than fitted constants.
\item \textbf{A validation suite} (\S\ref{sec:validation}) that replays a real
production trace --- the four Helios clusters~\cite{helios2021} --- and
reproduces the published FIFO-versus-shortest-job-first completion-time ratios to
a mean absolute error of 3.4\,\%, under three stated modelling choices we restate
wherever the result appears. It adds closed-form queueing checks, a list of
published numbers \texttt{fleetsim} deliberately does \emph{not} reproduce, and
one mechanism we had misdiagnosed (\S\ref{sec:replay-found}).
\end{enumerate}

\S\ref{sec:background} primes readers new to this infrastructure;
\S\ref{sec:models}--\ref{sec:traffic} give the model and traffic generator,
\S\ref{sec:validation} the validation, \S\ref{sec:cases}--\ref{sec:limitations}
the case studies, implementation, related work and limitations.

\section{Background: what a fleet is, and why scheduling it is hard}
\label{sec:background}

\subsection{The physical machine}

The unit of capacity is a \emph{chip}: one GPU or one TPU. Eight typically share
a server chassis and a fast internal switch --- a \emph{node}. Tens of nodes make
a rack or scalable unit, and tens of those a \emph{pod}: 4{,}096 GPUs in the H100
configurations we model, and 4{,}096 chips in a TPU v4 pod~\cite{tpuv4isca2023}.
Pods join into \emph{clusters}, then datacenters and \emph{metros}.

Every level is a bandwidth cliff, which is why locality matters. Llama~3's
cluster was eight pods of 3{,}072 GPUs at 1:7 oversubscription, and the 405B run
took 16{,}384 of those GPUs~\cite{llama3}. The 1:7 figure means seven times less
bandwidth between pods than within one, and a job straddling such a boundary
runs slower for its whole life.

\subsection{Gang scheduling and fragmentation}

A distributed training job is not a web service: its chips run one computation in
lockstep, synchronizing every step, so it needs \emph{all} of them \emph{at once}
for its whole life. Half a training job does not do half
the work; it does none. Allocation is therefore all-or-nothing --- \emph{gang
scheduling}. The analogy is a restaurant booking for twelve: a table for eight
plus a table for four across the street does not seat twelve.

The gang is also the unit of failure: if any node holding part of it dies, the
whole job dies. And gangs create \emph{fragmentation}, free capacity that cannot
be used: a fleet can have 200 chips free and still not start a 64-chip job,
because those chips sit in ones and twos across many nodes. Philly
measurements attribute 59--98\,\% of multi-GPU queueing delay to fragmentation
rather than fair-share policy~\cite{philly2019}; reproducing that attribution is
itself an explicit non-goal (\S\ref{sec:anti-goals}).

Three kinds of job share the machine, and they are not variations on a theme.
A \emph{pretraining} run trains a model from scratch: thousands of chips held in
lockstep for weeks. A \emph{finetune} adapts an already-trained model to a
narrower task: tens of chips for hours. An \emph{evaluation} scores an existing
model on a test set: one or two chips for minutes. All three arrive into the
same queue, and they differ by about four orders of magnitude in both size and
duration.

That disparity shows up in every trace. Over 82\,\% of Philly's jobs use one
GPU~\cite{philly2019}, and 70.8\,\% of jobs on the busiest Helios cluster we
replay are single-GPU. Acme measured evaluations at about 93\,\% of job
\emph{count} and under 1\,\% of GPU-\emph{time}, with pretraining at 1--3\,\% of
count and 70--94\,\% of GPU-time~\cite{acme2024}. The tiny jobs strand exactly
the capacity the enormous ones need.

This is the \emph{mice-and-hogs} shape: a swarm of one- and two-chip jobs (the
mice) that dominate the job count, and a handful of thousand-chip jobs (the
hogs) that dominate the chip-hours. Almost every metric in this paper has to be
reported twice because of it. Count one job as one job and you are describing
the mice; weight by chip-hours and you are describing the hogs.

A fleet is also shared. Each \emph{tenant} --- a team or project with its own
budget --- is granted a \emph{quota}, a promised share of the chips the
scheduler is expected to honour before it hands anything to anyone else. Quota
is a promise about capacity, not a physical partition: the chips themselves stay
in one pool.

\subsection{Three numbers that are not the same number}
\label{sec:threenumbers}

One term first, because the third definition below depends on it. A long
training job periodically writes its entire state to storage --- a
\emph{checkpoint} --- so that a job killed by a failure or a preemption resumes
from the last saved state rather than from the beginning. Saving costs the job
real time (about a minute, on \texttt{fleetsim}'s default one-hour interval),
and everything computed since the last checkpoint is lost when the job dies. The
analogy is a document editor that autosaves every hour: a crash costs you at
most an hour, and the autosaves themselves cost you a little throughput.

Conflating utilization metrics is the classic error here, so \texttt{fleetsim}
defines three and always reports all three. \emph{Allocation rate} is allocated
chip-time over \emph{total} chip-time, including chips that are broken or
drained. \emph{Occupancy} is allocated chip-time over \emph{healthy} chip-time.
\emph{Goodput} is the fraction of allocated chip-time that produced surviving
work, excluding work lost back to the last checkpoint, restart latency and
checkpoint-save overhead.

A worked example. Of 1{,}000 chips, fifty are down for repair, so 950 are
healthy, and all 950 are allocated. Allocation rate is $950/1000 = 0.95$;
occupancy is $950/950 = 1.00$, and the gap is exactly the capacity lost to
failures --- which is why occupancy alone makes a fleet with broken hardware look
perfect. Now suppose a third of the allocated chip-time goes on redoing work lost
at the last interruption and on restarts: goodput is $0.67$, and the share of the
machine doing useful work is $0.95 \times 0.67 \approx 0.64$. Fully occupied, a
third of its work wasted. Philly clusters were in fact measured at roughly
100\,\% allocation with about 52\,\% hardware utilization~\cite{philly2019}.

\subsection{Failures and preemption}
\label{sec:failures}

At small scale hardware failure is a nuisance; at fleet scale it is a design
constraint, because a gang fails when \emph{any} of its nodes fails. A gang of
$N$ nodes therefore breaks $N$ times as often as one node does.

Two published anchors bracket the numbers. Meta measured node mean time between
failures (MTBF) at about 42 node-days~\cite{metareliability2024}: one node in
the fleet breaks about every 42 days. Llama~3's 405B run, 16{,}384 GPUs and so
about 2{,}048 nodes, saw roughly one unplanned interruption every 3.1
hours~\cite{llama3}. Scaling that measurement, a gang eight times larger ---
16{,}384 nodes, the size our frontier case study uses --- would be interrupted
eight times as often, about every 23 minutes.

The two anchors do not agree with each other, and we say so rather than picking
the convenient one. Applying the 42-node-day MTBF directly to Llama~3's 2{,}048
nodes predicts an interruption every 30 minutes, roughly six times more often
than Llama~3 actually reported. The MTBF checks out against Meta's own smaller
figure (42 node-days over the 128 nodes of a 1{,}024-GPU job gives 7.9 hours,
which is what Meta measured) and not against Llama~3's. \texttt{fleetsim}
therefore ships 42 node-days as a configurable default, not as a fitted
constant, and the 23-minute figure quoted later comes from scaling Llama~3's
measurement, not from the MTBF.

Preemption --- stopping a running job to give its chips to a more important one
--- costs the same way: a preempted training job does not pause, it loses
everything since its last checkpoint and pays about 15 minutes of restart
overhead~\cite{metareliability2024}. Under a one-hour checkpoint interval that is
up to an hour of work per chip in the gang. This is how aggressive
preemption raises occupancy while destroying goodput.

\section{What fleetsim models}
\label{sec:models}

\subsection{One tree for two ecosystems}

Capacity is a tree of \emph{domains} from metro down to node, all one
\texttt{Domain} class. Leaves are nodes carrying a chip count, a chip type, a
health state (\texttt{HEALTHY}, \texttt{DRAINING}, \texttt{FAILED},
\texttt{MAINTENANCE}) and a failure-rate multiplier. Interior domains carry chip
counters maintained incrementally, making ``how much free capacity is under this
domain'' an $O(1)$ query.

One tree covers both ecosystems because \emph{level names are configuration, not
code}: a cluster declares \texttt{[cluster, pod, su, node]} for an H100 SuperPOD
or \texttt{[cluster, pod, cube, host]} for a TPU v5p pod. The scheduler sees
domains, capacities and constraints, never vendors.

\subsection{The allocation rule, and why it is the hardest call}

Jobs are allocated chip-granularly \emph{within} a single node and in whole-node
multiples \emph{above} one node: a 2-chip evaluation may share a node, a 64-chip
job asks for eight entire nodes. The reason is workload shape ---
\texttt{fleetsim}'s shipped class mix puts about 90\,\% of arrivals in 1--8-chip
evaluations, following the Acme characterization~\cite{acme2024}, which measured
the corresponding share at about 93\,\%. Were the node atomic,
a 1-chip evaluation would strand seven GPUs and corrupt both occupancy and the
mice-and-hogs story; rounding those jobs up would falsify the workload.

The rule has one load-bearing consequence, the mechanism behind our main
validation finding: \emph{a leaf with any owner at all is invisible to every
whole-node request}. A node with one chip in use has seven free chips no
multi-node gang can claim, so free capacity accumulates as 1--7-chip remainders.
That whole-node rule is an approximation we measured: in the Helios September
window the
trace's own node count satisfies $\texttt{node\_num} = \lceil
\texttt{gpu\_num}/8 \rceil$ for 99.66\,\% of multi-GPU jobs, and all 225
exceptions used \emph{more} nodes than the ceiling, never fewer --- the half the
model depends on.

An allocation is atomic, immutable once granted, and the unit of failure. A gang
may name a level it must fit \texttt{within}, and may be \emph{segmented} into
whole-node blocks that each sit inside one domain at that level.

\subsection{The engine}

Time is an \texttt{int64} count of microseconds, which makes tie-breaking exact.
The engine is event-driven rather than tick-driven because
state changes are sparse and unevenly spaced: a fleet can go minutes without one,
then process hundreds of preemptions at a single instant.

Events sit in a heap keyed by \texttt{(time, type, seq)}. The event type is not
decoration --- its integer value \emph{is} the same-timestamp ordering rank,
arranged so every repair, completion, failure and preemption at time $t$ is
processed before the scheduler wakes at $t$, so it always sees settled state.
Scheduler wakes are then
\emph{coalesced}: events set a dirty flag and ensure at most one pending wake at
the next round boundary, so the scheduler runs once per round (60\,s by default)
rather than once per event.

\subsection{Schedulers and placement}

Ordering and placement are separate axes~\cite{blox2024}: a scheduler sees an
immutable view of the cluster and returns \texttt{Place} and \texttt{Preempt}
intents, which the engine validates and applies~\cite{batsim2016}.

Four schedulers ship. \texttt{fifo} serves the queue in arrival order, with
strict head-of-line blocking by default: the head of the queue holds up
everything behind it until it can run. \texttt{tiered\_priority} implements
Borg-style priority bands~\cite{borg2015,borgtng2020} with requeue preemption.
\texttt{sjf} runs the shortest job first. \texttt{easy\_backfill} adds
\emph{backfilling}~\cite{easybackfill1995}: letting a small job at the back of
the queue jump ahead into a hole the head job cannot use, on the condition that
doing so provably does not delay the head. That condition is computed from each
job's declared runtime estimate, and users routinely over-estimate, so a
backfill scheduler must behave sanely when the estimate is wrong. Plugins
register through a Python entry point.

Orthogonally, four placement policies choose \emph{which} fitting nodes to take.
\texttt{first\_fit} (the default) takes the first nodes that fit, scanning by
ascending node identifier. \texttt{best\_fit} takes the nodes with the least
free space that still fits, so partly-used nodes fill up before empty ones are
opened --- \emph{tightest-fit} packing, a term used throughout the rest of the
paper. \texttt{consolidate} does the same at the sub-node level and additionally
minimises how many parent domains a multi-node gang touches; on a single-level
fleet it is bit-identical to \texttt{best\_fit}. \texttt{spread} takes the
emptiest nodes (worst fit) --- a deliberate anti-policy, our control arm.

\subsection{Failures, repair and maintenance}

Failures are sampled by \emph{aggregate} rate --- one exponential draw at the
fleet-wide rate, then a victim node --- not by per-node timers, so a job's mean
time to interruption scales inversely with its size for free.

Defaults are sourced, not invented: node MTBF 42 node-days, automatic repair
uniform on 60--180 minutes, 10\,\% of failures needing manual repair over 1--3
days, one maintenance drain per node-month, a one-hour drain grace window, and a
failure-cause mix of $\sim$60\,\% GPU/HBM, $\sim$10\,\% network and
$\sim$13\,\% software, with the remaining 17\,\% a residual ``other'' bucket
\texttt{fleetsim} adds so the mix sums to one. The MTBF comes from Meta's
measurements~\cite{metareliability2024}, the three approximate cause shares from
Llama~3~\cite{llama3}, and the drain rate from Borg~\cite{borg2015}. The repair
distributions are the weakest-sourced of the set: the repository attributes them
to AIReSim~\cite{aireesim}, a reliability simulator it names but does not pin to
a citable artefact, so a reader cannot check them the way the others can. All
are configuration values, none universal. Maintenance is a separate stream: a
\texttt{DRAINING} node takes no new placements, and its residents are requeued
when the grace window expires.

\subsection{Metrics}

Table~\ref{tab:metrics} states the definitions. State metrics are time-weighted
integrals updated on every state change --- exact, not sampled --- and
accumulated as integer chip-microseconds; discrete outcomes are event-sourced.

\begin{table}[t]
\centering
\caption{Metric definitions. ``Allocated'' means chips held by a running job;
``healthy'' excludes nodes that are failed, draining or under maintenance.
Integrals are exact integer chip-microseconds internally, reported as
chip-seconds.}
\label{tab:metrics}
\small
\begin{tabular}{@{}lp{0.36\linewidth}p{0.30\linewidth}@{}}
\toprule
Metric & Definition & What it tells you \\
\midrule
Allocation rate & $\int \text{allocated}\,dt \,/\, \int \text{total}\,dt$ &
Share of the purchased fleet handed out \\
Occupancy & $\int \text{allocated}\,dt \,/\, \int \text{healthy}\,dt$ &
Occupancy minus allocation rate is capacity lost to failures and drains \\
Goodput & $\sum \text{productive chip-s} \,/\, \int \text{allocated}\,dt$ &
Excludes work lost since the last checkpoint, restart latency and
checkpoint-save overhead \\
ETTR (per job) & Effective training time ratio: productive seconds / wall
seconds holding an allocation &
Goodput seen from one job's side; counts restart overhead and preemption grace
in the denominator \\
Fragmentation index & $1 - \text{largest placeable} / \text{total free}$ at a
level & $0$ when all free capacity is usable by one job \\
Stranded whole nodes & healthy leaves with $0 < \text{free} < \text{chips}$ &
Free capacity no whole-node gang can claim \\
\bottomrule
\end{tabular}
\end{table}

Two conventions matter. Metrics are given over the full run and over a
steady-state window (the middle 80\,\% of the horizon by default), and every
distributional metric is reported job-weighted \emph{and} chip-hour-weighted ---
a job-weighted mean over a mice-dominated workload describes the mice.
Best-effort jobs are excluded from wait distributions. \emph{Best-effort} is the
lowest, always-preemptible priority band, and \texttt{fleetsim} submits it from
a \emph{closed loop}: a source that keeps a fixed number of jobs in flight and
launches a replacement the instant one finishes. A closed loop by construction
always has work queued, so its queue length says nothing about the fleet and its
mean waiting time is undefined. Best-effort work is therefore reported by its
goodput instead.
